\documentclass[12pt,a4paper]{article}
\usepackage[utf8]{vietnam}
\usepackage{amsmath,amssymb}
\usepackage{geometry}
\geometry{
  top=12mm,
  bottom=12mm,
  left=14mm,
  right=14mm,
  footskip=6mm,
  headheight=14pt
}

\usepackage{tikz}
\usetikzlibrary{arrows.meta,calc,shapes.geometric}
\usepackage[most]{tcolorbox}
\usepackage{fontawesome}
\usepackage{fancyhdr}
\usepackage{enumitem}

\pagestyle{fancy}
\fancyhf{}
\fancyhead[L]{\footnotesize\color{teal!80!black}\faBook\ \textbf{SỔ TAY KIẾN THỨC TOÁN 3}}
\fancyhead[R]{\footnotesize\color{teal!80!black}\textbf{Lớp Toán Thầy Lĩnh -- loptoanthaylinh.vn}}
\fancyfoot[C]{\footnotesize\color{gray}-- Sổ tay kiến thức: Bài 07. Giảm đi một số lần --}
\renewcommand{\headrulewidth}{0.4pt}
\renewcommand{\footrulewidth}{0pt}

\setlength{\parindent}{0pt}
\setlength{\parskip}{0pt}

\begin{document}

% --- BANNER BÀI HỌC ---
\begin{tcolorbox}[
  enhanced,
  colback=teal!8!white,
  colframe=teal!70!black,
  arc=3mm, boxrule=1.2pt,
  left=3mm, right=3mm, top=2mm, bottom=2mm,
  halign=center
]
  {\LARGE\bfseries\color{teal!90!black}\faBookmark\ BÀI 07: GIẢM ĐI MỘT SỐ LẦN}\\[1mm]
  {\itshape\footnotesize\color{black!75} Quy tắc chia số đó cho số lần, phân biệt bản chất 'giảm đi số lần' và 'ít hơn' trên 1 trang A4}
\end{tcolorbox}

\vspace{-2mm}

\noindent
\begin{minipage}[t]{0.49\textwidth}\vspace{0pt}
\begin{tcolorbox}[
  enhanced,
  colback=blue!4!white,
  colframe=blue!75!black,
  arc=2mm, boxrule=0.8pt,
  left=2mm, right=2mm, top=1.8mm, bottom=1.8mm,
  title={\footnotesize\bfseries\color{white}\faCut\ 1. Quy tắc Giảm đi một số lần}
]
{\footnotesize
\begin{itemize}\setlength{\itemsep}{1pt}\setlength{\leftskip}{-2.5mm}
  \item \textbf{Quy tắc vàng:} Muốn giảm một số đi nhiều lần, ta lấy \textbf{số đó chia cho số lần}:
  $$\text{Số mới} = \text{Số đã cho} : \text{Số lần}$$
  \item \textit{Ví dụ:} Đoạn thẳng dài $18\text{ cm}$, giảm đi $3$ lần thì được đoạn thẳng dài: $18 : 3 = 6\text{ cm}$.
  \item Giảm đi $2$ lần nghĩa là giảm đi một nửa.
\end{itemize}
}
\end{tcolorbox}
\end{minipage}
\hfill
\begin{minipage}[t]{0.49\textwidth}\vspace{0pt}
\begin{tcolorbox}[
  enhanced,
  colback=teal!4!white,
  colframe=teal!75!black,
  arc=2mm, boxrule=0.8pt,
  left=2mm, right=2mm, top=1.8mm, bottom=1.8mm,
  title={\footnotesize\bfseries\color{white}\faBalanceScale\ 2. Phân biệt 'Giảm lần' \& 'Ít hơn'}
]
{\footnotesize
\begin{itemize}\setlength{\itemsep}{1pt}\setlength{\leftskip}{-2.5mm}
  \item \textbf{GIẢM ĐI $n$ lần} $\implies$ Thực hiện phép \textbf{CHIA} $(:)$:
  $$a \text{ giảm } n \text{ lần} \implies a : n$$
  \item \textbf{ÍT HƠN / BỚT ĐI $n$ đơn vị} $\implies$ Thực hiện phép \textbf{TRỪ} $(-)$:
  $$a \text{ ít hơn } n \text{ đơn vị} \implies a - n$$
  \item \textit{Ví dụ đối chiếu với số 20 và 4:}
  \begin{itemize}\setlength{\itemsep}{0.5pt}\setlength{\leftskip}{-2.5mm}
    \item $20$ giảm đi $4$ lần: $20 : 4 = \mathbf{5}$.
    \item $20$ ít hơn $4$ đơn vị: $20 - 4 = \mathbf{16}$.
  \end{itemize}
\end{itemize}
}
\end{tcolorbox}
\end{minipage}

\vspace{-2.5mm}
% --- KHỐI 3: SƠ ĐỒ TƯ DUY TIKZ ---
\noindent
\begin{tcolorbox}[
  enhanced,
  colback=orange!4!white,
  colframe=orange!80!black,
  arc=2mm, boxrule=0.8pt,
  left=2mm, right=2mm, top=1.5mm, bottom=1.5mm,
  title={\footnotesize\bfseries\color{white}\faSitemap\ SƠ ĐỒ TƯ DUY: MÔ HÌNH CHIA ĐOẠN THẲNG GIẢM 4 LẦN}
]
\centering
\begin{tikzpicture}[
  seg/.style={draw=blue!80!black, fill=blue!15, rounded corners=1pt, minimum height=0.45cm, font=\scriptsize\bfseries}
]
  \node[anchor=east, font=\scriptsize\bfseries] at (0, 0.7) {Đoạn ban đầu ($20\text{ cm}$):};
  \node[seg, minimum width=1.3cm] at (1.1, 0.7) {};
  \node[seg, minimum width=1.3cm] at (2.5, 0.7) {};
  \node[seg, minimum width=1.3cm] at (3.9, 0.7) {};
  \node[seg, minimum width=1.3cm] at (5.3, 0.7) {};
  \node[font=\scriptsize\bfseries, blue!80!black] at (3.2, 1.15) {Chia thành 4 phần bằng nhau};

  \node[anchor=east, font=\scriptsize\bfseries] at (0, 0) {Giảm 4 lần (1 phần):};
  \node[seg, minimum width=1.3cm, fill=teal!25, draw=teal!80!black] at (1.1, 0) {$5\text{ cm}$};

  \draw[thick, ->, purple!80!black] (1.1, 0.45) -- (1.1, 0.25) node[midway, right, font=\tiny\bfseries] {Lấy 1 phần: $20 : 4 = 5$};
\end{tikzpicture}
\end{tcolorbox}

\vspace{-2.5mm}
% --- KHỐI 4: VÍ DỤ & MẸO TRÁNH BẪY ---
\noindent
\begin{minipage}[t]{0.49\textwidth}\vspace{0pt}
\begin{tcolorbox}[
  enhanced,
  colback=green!4!white,
  colframe=green!60!black,
  arc=2mm, boxrule=0.8pt,
  left=2mm, right=2mm, top=1.8mm, bottom=1.8mm,
  title={\footnotesize\bfseries\color{white}\faCheckCircle\ Ví dụ thực hành chuẩn mực}
]
{\footnotesize
\textbf{Bài toán:}\\
Bao gạo to nặng $45\text{ kg}$, bao gạo nhỏ cân nặng giảm đi $5$ lần so với bao gạo to. Hỏi bao gạo nhỏ nặng bao nhiêu ki-lô-gam?\\
\textbf{Bài giải:}\\
Bao gạo nhỏ nặng số ki-lô-gam là:
$$45 : 5 = \mathbf{9}\text{ (kg)}$$
\textbf{Đáp số:} $9\text{ kg}$.
}
\end{tcolorbox}
\end{minipage}
\hfill
\begin{minipage}[t]{0.49\textwidth}\vspace{0pt}
\begin{tcolorbox}[
  enhanced,
  colback=red!4!white,
  colframe=red!70!black,
  arc=2mm, boxrule=0.8pt,
  left=2mm, right=2mm, top=1.8mm, bottom=1.8mm,
  title={\footnotesize\bfseries\color{white}\faExclamationTriangle\ Mẹo tránh bẫy trắc nghiệm}
]
{\footnotesize
\begin{itemize}\setlength{\itemsep}{1.5pt}\setlength{\leftskip}{-2.5mm}
  \item \textbf{Bẫy từ "giảm" làm phép trừ:}\\
  Học sinh rất hay nhầm: "Giảm đi 5 lần" làm thành $45 - 5 = 40$ (SAI!). Khi có từ \textbf{"lần"} bắt buộc phải làm phép \textbf{chia}.
  \item \textbf{Mẹo ghi nhớ đối xứng:}\\
  $\bullet$ Gấp ... lần $\to$ Nhân $(\times)$.\\
  $\bullet$ Giảm ... lần $\to$ Chia $(:)$.
\end{itemize}
}
\end{tcolorbox}
\end{minipage}


\end{document}
